% =============================================================================
% Consideracoes finais do Relatorio 2 (projeto conceitual do sistema Tag&Go).
% Sintese dos itens A a F com os numeros-chave, limitacoes e proximos passos.
% Sem figuras nem tabelas; os numeros sao os das Secoes 1 a 8 (biblia, 2.11).
% =============================================================================

\section{CONSIDERAÇÕES FINAIS}
\label{sec:r2-consideracoes}

No Relatório 2, a equipe percorreu os seis itens do projeto conceitual previstos no roteiro da disciplina \cite{zancul2026conceitual} e transformou a proposta do Relatório 1, uma etiqueta antifurto que o próprio cliente libera depois de pagar pelo celular, no conceito de um sistema completo, o Tag\&Go. O sistema reúne uma etiqueta ativa e reutilizável, a Plataforma Tag\&Go e seus serviços, dois caminhos de uso sobre um núcleo comum, uma infraestrutura de loja que atende às exceções e um ciclo logístico que devolve as etiquetas à fábrica ou ao centro de distribuição, onde elas são inspecionadas, recarregadas e cadastradas em novas peças. A Compra Rápida dispensa a instalação de aplicativo: o cliente abre o App Clip Tag\&Go no iPhone ou a Página de Compra Rápida no navegador do Android e paga com Apple Pay, Google Pay, cartão ou Pix copia e cola. A Jornada Completa roda no módulo Tag\&Go do app Riachuelo e acrescenta perfil, Ria Club, recomendação restrita ao estoque da loja no tamanho do perfil, compra on-line do tamanho ausente, com retirada na loja ou entrega sem frete acima de um valor mínimo, e métricas de sustentabilidade estimadas por categoria. Antes dos seis itens, a equipe corrigiu o único ponto em que o Relatório 1 perdeu nota e reescreveu as especificações-meta em três famílias, com sete requisitos do cliente (R01 a R07), doze especificações da etiqueta (E01 a E12) e dez especificações do sistema e do serviço (S01 a S10), cada uma com grandeza, unidade, valor-meta com tolerância e método de verificação (Tabela~\ref{tab:r2-especificacoes}).

\subsection{Síntese dos resultados}

Na análise funcional (Seção~\ref{sec:r2-funcional}), a equipe definiu a função global como liberar peça mediante pagamento (F0) e a representou numa caixa-preta com os fluxos de energia, material e sinal que atravessam a fronteira do sistema. Das duas estruturas funcionais comparadas, a liberação na própria peça (EF-1) foi reprovada em R07, por não oferecer caminho quando a etiqueta ou o celular falham, e a liberação em ponto fixo (EF-2) ficou fraca nos requisitos que motivaram o projeto, o tempo de pagamento (R03) e a compra sem funcionário (R04). A equipe escolheu, por isso, a estrutura mista EF-M, com liberação na peça no fluxo principal e liberação assistida no Ponto de Apoio para todas as falhas, e a desdobrou em 31 funções além da global: seis básicas (F1 a F6), 19 secundárias necessárias, três acessórias e três de estima. As seis básicas pertencem ao núcleo comum, o que confirma que os dois caminhos diferem apenas em funções secundárias. A localização do produto na loja (F25) foi mantida como secundária acessória, removível sem afetar F0. Entre os quatro níveis de granularidade avaliados, da unidade (N0) à peça (N3), a equipe adotou o N1, o setor e o piso, obtido do estoque RFID e de um planograma digital, sem depender da etiqueta; quando a leitura RFID tem mais de 24 horas ou apresenta divergência, o app indica a disponibilidade como ``provável'' e oferece a compra on-line. O N2 ficou como evolução, caso a Riachuelo instale leitores de teto, e a etiqueta como farol de rádio foi rejeitada porque, pela estimativa da equipe, esgotaria a célula em cerca de 16 dias, menos que um ciclo. As funções do ciclo da etiqueta (F17 a F24) registram a mudança mais profunda em relação ao Relatório 1: a etiqueta deixou de ser consumível e passou a ser um ativo retornável.

No estudo de diferenciação (Seção~\ref{sec:r2-diferenciacao}), a equipe comparou o Tag\&Go com dez referências em sete atributos, da rede tradicional e das lojas conceito da Riachuelo a soluções de outros mercados, como MishiPay, Paytag, rapitag e Sensormatic SuperTag 4. No mapa de posicionamento, construído com posições estimadas pela equipe, o quadrante que combina liberação sem ponto fixo e barreira autenticada na peça ficou com apenas duas soluções: a rapitag, que exige o app do varejista, e o Tag\&Go. A equipe formulou cinco diferenciais (DF1 a DF5), cada um pelo contraste com uma referência concreta: pagamento em qualquer ponto da loja com liberação na própria peça; autorização verificada na etiqueta, com imunidade ao destacador magnético; Compra Rápida sem instalar aplicativo nas duas plataformas; ciclo fechado de uma etiqueta ativa com cadastro na origem; e Jornada Completa ligada ao Ria Club. Registrou também o que não diferencia o produto, como o self-checkout em si, o RFID, a carteira digital e o cashback, que o mercado já oferece. Os riscos de diferenciação são a entrada da Paytag ou da rapitag no Brasil e a tecnologia não divulgada das lojas conceito da Riachuelo, que pode já cobrir parte do DF1.

Na escala vertical e no valor mercadológico (Seção~\ref{sec:r2-valor}), a equipe tratou o valor como uma questão de viabilidade para dois agentes, a Riachuelo e o fornecedor, porque o Tag\&Go é vendido a um varejista e não tem concorrente direto à venda no Brasil. A escala de 25 itens, com 23 produtos de referência em cinco categorias, posicionou o custo-meta de R\$~60 entre as etiquetas eletrônicas de prateleira (R\$~26,00 a R\$~51,99) e os rastreadores BLE de consumo (R\$~129,93 a R\$~369,00), posição coerente com um dispositivo que tem rádio, bateria e atuador. Como a etiqueta custa cerca de 70 vezes a mini tag que substitui, a equipe concluiu que ela só se viabiliza como ativo reutilizável, cobrado por uso. O custo-meta foi fixado em R\$~60 $\pm$ 10\%, com composição de R\$~63,25, contra R\$~122,00 na escala sem redesenho e R\$~136,21 no protótipo; a redução de 48,2\% veio de três decisões de projeto, a célula menor, o micromotor com came e a placa única. A frota foi dimensionada, como premissa, em 6.854 etiquetas por loja ($68 \times 84 \times 1{,}2$). Na locação, modalidade principal, a Riachuelo paga R\$~2,00 por etiqueta por mês, R\$~800 por mês pela plataforma e R\$~15.000 de implantação por loja, o que equivale a R\$~7,01 por liberação, contra R\$~0,29 com a etiqueta rígida atual, e empata ao evitar 3,68 desistências de compra por dia por loja, 1,6\% das 227 compras diárias estimadas. Para o fornecedor, esse preço dá TIR de 22,0\% em 5 anos, acima da taxa mínima de atratividade de 20\% adotada como premissa; a R\$~1,80, a TIR cairia para 16,6\%, e o preço mínimo para remunerar o capital é de R\$~1,92. Na análise de sensibilidade, a equipe verificou que a fração de peças com etiqueta, a taxa de retorno e as perdas pesam mais no equilíbrio que o próprio preço.

No estudo de aproveitamento técnico (Seção~\ref{sec:r2-aproveitamento}), a equipe organizou a busca em oito linhas de similaridade, de L1 a L8, e verificou que quase todas as funções do Tag\&Go têm precedente comercial isolado: a retenção vem das etiquetas rígidas, a leitura e o pagamento sem aplicativo vêm das plataformas móveis, a recirculação vem da Sensormatic e da Pact, e a localização por setor vem da Zara e do Walmart. O desenvolvimento próprio ficou concentrado na integração dessas soluções e em dois pontos sem solução pronta: a verificação do token na própria etiqueta, com desafio gerado por ela, e a recirculação de uma etiqueta com bateria e chave criptográfica. As soluções descartadas, como a visão computacional, o alarme sem barreira física, o pagamento facial próprio e o anúncio BLE contínuo, ficaram registradas com o motivo. Na propriedade intelectual, a patente mais próxima, a US 10.121.339 B2, está ativa até 2035 e tem uma reivindicação que reúne os elementos centrais da concepção escolhida, o celular, a compra verificada, o comando sem fio e o pino liberado sem ajuda humana, e nenhuma das famílias ativas levantadas tem publicação brasileira no Google Patents. A equipe registrou as diferenças técnicas do Tag\&Go em relação a essa patente sem tratá-las como parecer de liberdade de operação.

Na reformulação dos desenhos, a primeira parte (Seção~\ref{sec:r2-concepcao}) levantou os efeitos físicos e os portadores de cada função, combinou-os numa matriz morfológica e formou três alternativas, que passaram por um filtro eliminatório antes da matriz de decisão. Com dez critérios em dois blocos, o do cliente com 75\% do peso e o do varejista com 25\%, a A1 Etiqueta Ativa obteve 4,06 pontos, contra 3,85 da A3 Etiqueta Passiva Energizada e 3,72 da A2 Estação Inteligente. A escolha da A1 é condicional em dois sentidos. A A2 passa à frente quando o bloco do varejista pesa 36,6\% ou mais, porque a A1 custa R\$~414.458 por loja, 14,4 vezes o custo da A2 (R\$~28.877); e o fluxo do iPhone depende de o Core Bluetooth funcionar dentro de um App Clip, o que a equipe deduziu da documentação da Apple e ainda precisa testar. A A2 ficou registrada como concepção de reserva e como etapa intermediária de implantação, e a A3, como evolução futura sem bateria.

A segunda parte (Seção~\ref{sec:r2-arquitetura}) desdobrou a A1 numa arquitetura modular com cinco subsistemas, a etiqueta, a Plataforma, os canais do cliente, a infraestrutura de loja e o ciclo logístico, e dezoito interfaces classificadas por tipo de fluxo. O desenho da etiqueta manteve o envelope de $98 \times 82 \times 26$\,mm e mudou o arranjo interno: os contatos de serviço foram para a face superior, o servo deu lugar a um micromotor com came, a célula passou a 150\,mAh com o modo botão, que eleva a autonomia estimada a 448 dias contra a meta de 168, a sinalização ganhou som e padrões de piscada, a identificação ganhou NFC e o carretel não ferromagnético foi adotado. A segurança se apoia num token HMAC-SHA256 com chave por etiqueta, ligado a um nonce gerado pela própria etiqueta, válido por até 120\,s e de uso único, de modo que o celular funciona apenas como retransmissor. O ciclo da etiqueta, com ciclo médio de 84 dias e dois ramos de cadastro na origem, a fábrica de Natal e os CDs, tira da loja o trabalho de etiquetar sem eliminar o seu custo, que passa a R\$~0,14 por peça na origem.

Na comercialização e na distribuição (Seção~\ref{sec:r2-comercializacao}), a equipe tratou o Tag\&Go como uma venda entre empresas apoiada numa experiência de consumo. A venda à Riachuelo será direta e consultiva, começando por um piloto numa loja conceito, limitado a uma categoria, e seguindo para um contrato-quadro com as 20 lojas do SOM, que exigiriam 137.080 etiquetas e gerariam ao fornecedor R\$~3,48 milhões por ano. Os parceiros de integração que já atendem a Riachuelo entram como canais de implantação, e a expansão para outras redes do segmento aproveita a arquitetura modular, com a troca dos módulos de integração. A distribuição física segue o caminhão da frota própria na ida e no retorno, com recondicionamento em Natal e em Guarulhos. A equipe escolheu a locação como modalidade principal, com franquia de perda de 1\% das liberações, porque ela transfere ao fornecedor o risco de perda e de obsolescência da frota e sai mais barata para a Riachuelo quando se considera o custo de capital da compra. O tratamento de dados pessoais foi organizado por finalidade, com o mínimo necessário na Compra Rápida e consentimento separado por finalidade na Jornada Completa.

Em conjunto, a equipe concluiu que o Tag\&Go é tecnicamente coerente com as especificações-meta e economicamente plausível para a Riachuelo sob três condições, que precisam ser verificadas antes de qualquer decisão de implantação: evitar pelo menos 3,7 desistências de compra por dia por loja, manter o retorno das etiquetas perto da meta de 99\% (S01), acima da faixa de 80\% a 83\% observada na recirculação de etiquetas rígidas, e restringir o uso às peças que hoje já recebem etiqueta rígida, de preferência as de maior valor e de giro rápido. A escolha da concepção depende, além disso, do peso que a Riachuelo atribuir ao custo e à operação e do teste de Core Bluetooth no App Clip.

\subsection{Limitações}

A principal limitação do estudo é que a economia do sistema repousa sobre premissas ainda não medidas. A fração de peças com etiqueta rígida ($f = 10\%$), a permanência de 60 dias da peça na loja, que define o ciclo médio de 84 dias, o ticket de R\$~210, a adesão de 24\% das compras pelo celular e o próprio custo-meta foram estimados pela equipe a partir de fontes públicas, e a permanência em loja foi ancorada no prazo médio de estoques da Renner porque o indicador da Riachuelo não foi verificado. O número de desistências evitadas, variável que decide o equilíbrio, não tem estimativa pública para a rede. Na análise de sensibilidade, variações plausíveis dessas premissas levaram o equilíbrio de 1,82 a 7,40 desistências evitadas por dia por loja, no caso da fração de peças etiquetadas, e a 9,16, se o retorno cair ao nível do mercado.

A segunda limitação está na voz do comprador. O QFD do Relatório 1 foi construído com a voz do consumidor, e os critérios do varejista na matriz de decisão (K7 a K10), assim como a divisão de 75\% e 25\% entre os blocos, são premissas da equipe. Não houve pesquisa com o comprador B2B, ao contrário de precedentes da disciplina que posicionaram o produto na escala vertical com pesquisa de clientes. Como a A2 passa à frente com um peso de 36,6\% para o bloco do varejista, a escolha da concepção depende de uma validação que ainda não ocorreu.

A terceira limitação diz respeito à qualidade dos preços. Parte dos preços de equipamentos de automação comercial veio de páginas e de títulos de anúncio do varejo brasileiro, com confiança média, e os preços em dólar, de distribuidores, sem impostos de importação, foram convertidos pela PTAX de um único dia, 25/09/2026. Os preços de volume da célula, do micromotor e do molde, que sustentam o custo-meta, ainda não foram cotados, e a carcaça injetada foi estimada a partir de um exemplo que não constitui orçamento.

A quarta limitação é o tratamento das perdas. A equipe fixou em zero o efeito do sistema sobre as perdas no caso base, porque as evidências apontam em sentidos opostos. No cenário pessimista, baseado no estudo do ECR de 2026, em que 35 dos 39 varejistas analisados eram supermercados, o equilíbrio sobe para 6,6 a 8,3 desistências evitadas por dia por loja, e a transposição desse resultado para o varejo de moda não foi verificada.

A quinta limitação envolve as plataformas móveis. O funcionamento do Core Bluetooth dentro de um App Clip é uma dedução da documentação e de uma resposta da Apple no fórum oficial, sem confirmação explícita, e o Pix pelo Google Pay com Open Finance na web é hipótese a validar com o PSP; nenhum dos dois foi testado. Na propriedade intelectual, a equipe consultou apenas o Google Patents, e a ausência de publicação brasileira das famílias levantadas ainda precisa ser confirmada no INPI. Por fim, nenhuma especificação da etiqueta foi medida em bancada: a margem de força do atuador, o tempo de liberação, a autonomia e as liberações por carga são estimativas calculadas a partir de dados de componentes e dos estudos anteriores da equipe, e a meta de resistência à extração do pino ainda precisa ser comparada com a de uma etiqueta comercial.

\subsection{Próximos passos}

Os Relatórios 3 e 4 e o protótipo obrigatório da entrega final deverão converter essas premissas e estimativas em medições. A equipe definiu sete frentes de trabalho:
\begin{itemize}
    \item cotar em volume a célula, o micromotor e o molde das carcaças e detalhar formas, materiais e processos de fabricação, com o projeto para manufatura e a lista de materiais detalhada, para confirmar o custo-meta de R\$~60 $\pm$ 10\%;
    \item construir o protótipo com o App Clip Tag\&Go e a Página de Compra Rápida e executar o primeiro marco de go/no-go, que verifica se o Core Bluetooth funciona num App Clip real; se o teste falhar, o iPhone continuará pagando pelo App Clip, e a liberação passará ao Liberador de Balcão;
    \item ensaiar as especificações da etiqueta, de E01 a E12, e a liberação condicionada ao pagamento (R06) com ataque, com meta de nenhuma liberação sem pagamento confirmado em pelo menos 200 tentativas;
    \item testar a usabilidade dos requisitos de tempo de pagamento, compra sem funcionário e clareza do fluxo (R03 a R05), com cronometragem por etapa do fluxo e amostra de pelo menos 20 pessoas;
    \item validar com a Riachuelo a fração de peças com etiqueta ($f$), o ciclo médio ($T$), o protocolo para estimar as desistências evitadas ($d$) e os pesos do bloco do varejista na matriz de decisão;
    \item elaborar, com o Jurídico da Riachuelo, o relatório de impacto à proteção de dados pessoais antes do piloto;
    \item fazer a busca no INPI das famílias de patentes levantadas no Google Patents.
\end{itemize}

A sequência dessas frentes segue o calendário da disciplina, com o Relatório 3 em 04/11, o Relatório 4 em 23/11 e a apresentação final com o protótipo em 07 e 09/12 \cite{zancul2026conceitual}. O teste de Core Bluetooth e a validação dos pesos vêm primeiro, porque cada um deles pode mudar a concepção escolhida; os ensaios e os testes de usabilidade vêm em seguida, porque dependem do protótipo. Depois da disciplina, o piloto numa loja conceito da Riachuelo deverá medir o retorno das etiquetas, o ciclo médio, os tempos por etapa e as desistências evitadas, e substituir por dados da operação as premissas que hoje sustentam a conta econômica. A equipe concluiu que, com essas medições, a Riachuelo poderá decidir a implantação com base no próprio desempenho do sistema, a tempo de alimentar a avaliação financeira do ciclo de reformas e aberturas de lojas, prevista para abril de 2027 \cite{economicnews2026riachuelo}.
